\chapter{Changing parameters}
\label{ch:params}

Parameters live in four places: the workbench forms, the workspace
\code{config.json}, the command-line flags of each service, and constants in
the code. This chapter says which is which, what each parameter does, and how
to change it safely.

\section{Where a parameter takes effect}

\begin{figure}[htbp]
\centering
\scalebox{0.8}{%
\begin{tikzpicture}[node distance=7mm and 12mm]
  \node[hot,minimum width=34mm] (ui) {workbench forms\\speed, data, model};
  \node[box,right=of ui,minimum width=34mm] (cfg) {\code{config.json}\\in the workspace};
  \node[box,right=of cfg,minimum width=34mm] (sup) {supervisor\\builds command lines};
  \node[box,below=of sup,minimum width=34mm] (svc) {services\\plant, PLC, historian, HMI};
  \node[box,below=of cfg,minimum width=34mm] (jobs) {jobs\\dataset, training};
  \node[bad,below=of ui,minimum width=34mm] (code) {code constants\\\code{Params}, \code{PidTuning}, tags};
  \draw[flow] (ui) -- node[lbl,above]{saves} (cfg);
  \draw[flow] (cfg) -- (sup);
  \draw[flow] (sup) -- node[lbl,right]{CLI flags} (svc);
  \draw[flow] (ui) |- ([yshift=2mm]jobs.west);
  \draw[flow] (code) -- node[lbl,below]{defaults} (jobs);
  \draw[flow] (code.south) -- ++(0,-6mm) -| node[lbl,pos=0.25,below]{defaults} (svc.south);
\end{tikzpicture}}
\caption{How a parameter reaches the code that uses it.}
\label{fig:paramflow}
\end{figure}

\begin{warn}[What is and is not wired]
The supervisor passes only \code{dt}, \code{speed}, \code{loop.setpoint} and
\code{loop.scan_period} from \code{config.json} to the services. The
\code{plant} section and \code{loop.kc}, \code{ti}, \code{td}, \code{tf} are
saved and read back, but the live services always use the code defaults of
\code{Params} and \code{PidTuning}. To experiment with plant parameters or
tuning, use the offline harness (Section~\ref{sec:offline}) or change the
defaults in the code.
\end{warn}

\section{The workspace configuration}

\code{config.json} is created on first start and rewritten atomically (to a
temporary file, then moved into place) whenever the workbench saves.

\begin{lstlisting}[style=sh]
{
  "name": "demo",
  "speed": 1.0,            // real-time factor of the live stack
  "dt": 0.05,              // s, plant integration step
  "ports": { "plant_modbus": 5400, "plant_control": 5401, "plc_modbus": 5402,
             "plc_opcua": 5403, "historian_api": 5404, "hmi": 5405 },
  "plant": { "area_tank1": 0.02, "area_tank2": 0.02, "height_max": 0.6,
             "pump_dp0": 60000.0, "pump_qmax": 0.000417,
             "valve_cv": 2.429e-06, "valve_rangeability": 30.0,
             "tau_valve": 2.0, "cda_12": 6.9e-05, "cda_2out": 9e-05 },
  "loop":  { "kc": 0.35, "ti": 72.0, "td": 0.0, "tf": 5.0,
             "setpoint": 300.0, "scan_period": 0.1 },
  "data":  { "duration": 3600.0, "magnitudes": 3, "repeats": 2,
             "normal_runs": 10, "window": 60, "stride": 10,
             "quantile": 0.995, "seed": 0 }
}
\end{lstlisting}
(JSON has no comments; they are added here for explanation.) Edit it with the
workbench stopped, or use the forms. A missing key falls back to its default.

\subsection{Runtime parameters}
\begin{center}
\begin{tabularx}{\textwidth}{@{}lllX@{}}
\toprule
Parameter & Default & Set from & Effect \\
\midrule
\code{speed} & 1.0 & \ui{Runtime} menu & real-time factor of plant and PLC (1, 2, 4, 10 or 20); saved when the stack starts \\
\code{dt} & 0.05\,s & \code{config.json} & RK4 integration step of the plant \\
\code{loop.setpoint} & 300\,mm & \code{config.json} & initial LIC101 setpoint; the operator can change it live \\
\code{loop.scan_period} & 0.10\,s & \code{config.json} & PLC scan period $T_s$ \\
\code{ports} & base 5400 & \code{--base-port} & six consecutive ports from the base \\
\bottomrule
\end{tabularx}
\end{center}

\begin{warn}[Speed on Windows]
At \code{speed} $s$ the PLC wakes every $T_s/s$ wall seconds. Windows timers
have a granularity of about 15.6\,ms, so above $4\times$ the scan period
($25$\,ms at $4\times$) approaches the timer resolution and scans overrun.
Use $\le 4\times$ on Windows; slow faults such as stiction need a few process
time constants (72\,s each) to develop.
\end{warn}

\subsection{Dataset parameters (Data tab)}
\begin{center}
\begin{tabular}{@{}lllll@{}}
\toprule
Parameter & Meaning & Quick & Standard & Thorough \\
\midrule
\code{duration} & seconds per run & 1200 & 3600 & 3600 \\
\code{magnitudes} & scale factors per fault & 1 & 3 & 5 \\
\code{repeats} & seeds per magnitude & 1 & 2 & 3 \\
\code{normal_runs} & minimum normal runs ($\ge3$) & 6 & 10 & 16 \\
\midrule
runs generated & & 14 & 58 & 136 \\
approximate time & & 1.5\,min & 4\,min & 10\,min \\
\bottomrule
\end{tabular}
\end{center}

Guidance: more magnitudes give the detector a wider view of fault severity;
more repeats make the evaluation numbers less noisy; more normal runs give a
better threshold and a more trustworthy false alarm rate. Duration must leave
time after onset (25--60\% of the run) for ramped faults to develop.

\subsection{Model parameters (Model tab)}
\begin{center}
\begin{tabularx}{\textwidth}{@{}llX@{}}
\toprule
Parameter & Default & Effect \\
\midrule
methods & all four & which detectors to fit and compare \\
\code{window} & 60 samples & longer windows see slower faults and smooth noise, at the cost of delay \\
\code{stride} & 10 samples & spacing of windows; smaller means finer delay resolution and more data \\
\code{quantile} & 0.995 & threshold quantile on calibration data; raise it (0.999) for fewer false alarms, lower it for faster detection \\
\code{seed} & 0 & split and model randomness (Python only) \\
\bottomrule
\end{tabularx}
\end{center}

\subsection{Fault parameters (Faults tab)}
\begin{center}
\begin{tabular}{@{}lll@{}}
\toprule
Kind & Starting magnitude & Unit \\
\midrule
\code{leak} & $5\times10^{-5}$ & m$^2$ effective leak area \\
\code{fouling} & 0.35 & fraction of flow area lost \\
\code{pump_wear} & 0.18 & fraction of head lost \\
\code{stiction} & 0.10 & deadband, fraction of stroke (slip default 0.04) \\
\code{sensor_bias} & 25 & tag units \\
\code{sensor_drift} & 1.5 & tag units per minute \\
\code{sensor_stuck} & none & \\
\code{sensor_noise} & 2.0 & $\sigma$ in tag units \\
\bottomrule
\end{tabular}
\end{center}
Sensor faults also take a tag (\tg{LT101}, \tg{LT102}, \tg{FT101},
\tg{FT102}, \tg{PT101} or \tg{ZT101}) and every fault takes a ramp in
seconds.

\section{Command-line options}
\label{sec:cli}

\subsection{Workbench}
\begin{lstlisting}[style=sh]
python -m pyplant.workbench.app [--workspace ./work/demo] [--base-port 5400]
                                [--host 127.0.0.1] [--port 8050] [--debug]
\end{lstlisting}

\subsection{Services}
These are the processes the supervisor starts; the \ui{Runtime} tab shows the
exact command line of each. They can be run by hand in four terminals.

\begin{center}
\begin{tabularx}{\textwidth}{@{}p{62mm}lX@{}}
\toprule
Flag & Default & Meaning \\
\midrule
\multicolumn{3}{@{}l}{\code{python -m pyplant.services.plant}} \\
\quad\code{--port} & 5020 & Modbus slave port \\
\quad\code{--control-port} & 0 & fault injection HTTP API; 0 disables \\
\quad\code{--dt} & 0.05 & integration step, s \\
\quad\code{--speed} & 1.0 & real-time factor \\
\quad\code{--start} & steady & \code{steady} or \code{empty} initial state \\
\quad\code{--no-noise} & off & disable instrument noise \\
\midrule
\multicolumn{3}{@{}l}{\code{python -m pyplant.services.plc}} \\
\quad\code{--plant-host}, \code{--plant-port} & 127.0.0.1, 5020 & where the plant is \\
\quad\code{--port} & 5021 & PLC Modbus server port \\
\quad\code{--opcua-port}, \code{--no-opcua} & 4840 & OPC UA server \\
\quad\code{--period} & 0.10 & scan period, s \\
\quad\code{--speed} & 1.0 & real-time factor \\
\quad\code{--setpoint} & 300 & initial setpoint, mm \\
\midrule
\multicolumn{3}{@{}l}{\code{python -m pyplant.services.historian}} \\
\quad\code{--endpoint} & & PLC OPC UA endpoint \\
\quad\code{--db} & pyplant.duckdb & database file \\
\quad\code{--publish-ms} & 200 & subscription publishing interval \\
\quad\code{--flush} & 1.0 & buffer flush period, s \\
\quad\code{--heartbeat} & 30 & maximum gap per tag, s \\
\quad\code{--api-port} & 0 & query API port; 0 disables \\
\midrule
\multicolumn{3}{@{}l}{\code{python -m pyplant.services.hmi}} \\
\quad\code{--endpoint} & & PLC OPC UA endpoint \\
\quad\code{--history} & & historian URL, or a DuckDB path \\
\quad\code{--port}, \code{--poll} & 8080, 0.5 & HTTP port, poll period (s) \\
\bottomrule
\end{tabularx}
\end{center}

\subsection{Training from the command line}
\begin{lstlisting}[style=sh]
python -m pyplant.ml.train --out artifacts --duration 3600 --magnitudes 3 \
       --repeats 2 --normal-runs 6 --window 60 --stride 10 --quantile 0.995 \
       --seed 0 [--reuse]
\end{lstlisting}
This fits PCA and the autoencoder; \code{--reuse} skips generation when
\code{dataset.parquet} already exists in \code{--out}.

\section{Code constants}

Some parameters are design decisions and live in code. Change them with care
and run \code{pytest} afterwards; the tests check most invariants.

\begin{center}
\begin{tabularx}{\textwidth}{@{}llX@{}}
\toprule
Constant & File & Meaning \\
\midrule
\code{Params} fields & \code{core/hybrid_tank.py} & plant geometry, pump, valve, orifices (Table~\ref{tab:params}) \\
\code{PidTuning} fields & \code{plc/pid.py} & $K_c$, $T_i$, $T_d$, $T_f$, limits, action \\
alarm limits & \code{tags.py} & \code{limits} on each tag: HI 500, HIHI 550, LOLO 50\,mm; PT HIHI 58\,kPa; FT LO 1\,L/min \\
deadbands & \code{plc/program.py} & 10\,mm, 2\,kPa, 0.5\,L/min \\
\code{MIN_FLOW_DELAY} & \code{plc/program.py} & 15\,s low-flow trip delay \\
\code{DEFAULT_NOISE} & \code{faults/scenarios.py} & instrument noise $\sigma$ per tag \\
\code{NORMAL_MOVES} & \code{faults/scenarios.py} & setpoint programme of the scenarios \\
\code{WINDOW}, \code{STRIDE} & \code{ml/features.py} & default window and stride \\
\code{PERSISTENCE} & \code{ml/evaluate.py} & consecutive windows for a detection (3) \\
\code{DEVELOPED_AFTER} & \code{ml/evaluate.py} & settling time before isolation is judged (300\,s) \\
\code{DATA_PRESETS}, \code{PRESETS} & \code{workbench/app.py} & dataset and fault presets in the UI \\
\bottomrule
\end{tabularx}
\end{center}

\section{Offline experiments with other parameters}
\label{sec:offline}

The fastest way to try a different plant or tuning is the offline harness: the
same PLC program, the same physics, no network, about 1900 times real time.

\begin{lstlisting}
from pyplant.core.hybrid_tank import Inputs, Params
from pyplant.plc.harness import OperatorPanel, PlcHarness, scheduled_setpoint
from pyplant.plc.pid import PidTuning
from pyplant.plc.program import Program
from pyplant.runners.base import Simulation
from pyplant.runners.fast import run

plant  = Params(area_tank1=0.03)                 # a wider tank 1
tuning = PidTuning(kc=0.6, ti=50.0)              # a more aggressive loop
harness = PlcHarness(program=Program(ts=0.1, tuning=tuning),
                     panel=OperatorPanel(setpoint=300.0))
sim = Simulation.at_steady_state(Inputs(0.662, True), params=plant,
                                 dt=0.05, control_period=0.1, policy=harness)
sim.policy = scheduled_setpoint(harness, [(0.0, 300.0), (60.0, 350.0)])
tr = run(sim, duration=600.0, log_period=1.0, extra=harness.tags)
df = tr.to_dataframe()            # LT101, LIC101_OP, ... one row per second
\end{lstlisting}

To generate scenario data with a different plant, pass \code{params} to the
scenario: \code{Scenario("wide", params=Params(area_tank1=0.03), ...)}. Note
that \code{run_scenario()} always uses the default \code{PidTuning}.

\subsection{Retuning after a plant change}
If you change the plant, re-identify $K$ and $\tau$ from an open-loop step
(Chapter~\ref{ch:tutorials}) and recompute the IMC settings:
\[
  K_c = \frac{\tau}{K\lambda}, \qquad T_i = \tau, \qquad
  \lambda \approx \tau/5 .
\]
A larger tank area increases $\tau$ roughly in proportion; the gain $K$
changes with the operating point because the valve and the drains are
non-linear.
